% ch10-summary.tex – Chapter 10: summary (final round, Writer). No new facts.
\ChapterMission{המשימה: לסכם את הממצאים, להעריך את מוכנות הפרויקט ולקבוע את הצעדים הבאים}
\ChapterIcon{pedestrian}
\chapter{סיכום}
\begin{scenebox}
ערב. על המדרכה עומדים ארבעה אנשים, ואחד מהם הגיע זה עתה. איש אינו יורד לכביש.
האיש הירוק נדלק, הקבוצה חוצה יחד, והספרות הכתומות מלוות אותה עד הצד השני.
\end{scenebox}

הקבוצה שחוצה יחד בירוק היא התוצאה שהפרויקט מבקש להשיג גם כשעל המדרכה עומד אדם אחד בלבד. הפרק מסכם את הממצאים ואת המוכנות, ומגדיר את הצעדים הבאים.

\section{ממצאים ומוכנות}
הולך רגל הממתין לבדו חוצה באדום יותר מהולך שממתין עם אחרים \cite{rosenbloom2009}, והנטייה להפר עולה ככל שההמתנה מתארכת (פרק~\ref{chap:problem}). הניתוח הביקורתי אישר את כיוון הממצא, אך הראה שגודל האפקט שונה מן המודפס (סעיף~\ref{sec:critical}).
הספירה בזמן האדום נדחתה בגלל התקן והעדות המעורבת, ובמקומה נבחרו אדום מסתגל בגבולות בטוחים, ספירה בפינוי וקמפיין (פרק~\ref{chap:alternatives}).
הרכיבים הטכניים קיימים כמוצרי מדף (פרק~\ref{chap:market}), לכל דרישה יש יעד ושיטת אימות (פרק~\ref{chap:requirements}), אף תמונה אינה נשמרת (פרק~\ref{chap:architecture}), ולתוכנית העבודה ולסיכונים הקריטיים יש אחראים ותוכניות חלופיות (פרקים~\ref{chap:workplan} ו-\ref{chap:risks}).
הפער העיקרי שנותר הוא רגולטורי: היתר לקיצור האדום והתאמת הגבולות הבטוחים, שנלקחו מתקנים אמריקניים, לתקן הישראלי.

\section{הצעדים הבאים}
ארבעה צעדים פותחים את הביצוע: פנייה לעירייה לבחירת אתר ולפתיחת הליך האישור; מיפוי התקן הישראלי וקביעת הגבולות הבטוחים לאתר; בדיקת מעבדה של שני גלאים לפחות מול דרישות הדיוק והפרטיות; ותכנון מדידת הבסיס, כדי לבחון את ההשערות בניסוי לפני ואחרי.
